\section{The engineering problem: a search returns hits}
\label{sec:postgres-custom-scan-scope}

The sequence interfaces in the previous chapter let an alignment algorithm
read a stored sequence without first copying all its bases. We now need to
connect that algorithm to SQL. The main difficulty is the result: one pair of
sequences can produce several alignments, and each alignment has several
fields. We call one returned alignment a \emph{hit}.

Consider a table \texttt{targets} with an identifier and a sequence. The SQL
examples in this chapter assume BioSeqDB's table access method and a DNA5
sequence column packed into 64-bit storage units:

\noindent\begin{minipage}{\linewidth}
\begin{lstlisting}[style=bioseq-search-sql,caption={Search once and return the fields of each hit.}]
CREATE TABLE targets (
  id bigint PRIMARY KEY,
  seq bioseq.seq('dna5', 64)
) USING bioseq;

SELECT t.id,
       bioseq.dp_score(t.seq) AS score,
       bioseq.target_start(t.seq) AS start,
       bioseq.cigar(t.seq) AS cigar
FROM targets AS t
WHERE t.id > 20
  AND mmsearch(t.seq, $1::bioseq.seq('dna5', 64));
\end{lstlisting}
\end{minipage}

The first \texttt{mmsearch} argument is the target sequence and the second
is the query. Here \texttt{\$1} supplies that query sequence. The
\texttt{score} alias means the dynamic-programming score of the current hit.
The CIGAR string describes which positions align and where insertions and
deletions occur.\footcite{sam_format_specification}
Three requirements follow from this example:

\begin{itemize}
  \item \textbf{Compute the search once.} Score, start, and CIGAR must read
    fields of the same hit. Calling the aligner separately for each field
    would repeat its most expensive work.
  \item \textbf{Keep ordinary SQL conditions useful.} The identifier test
    should remove unwanted targets before alignment. PostgreSQL should still
    see the tables, their types, and the surrounding joins.
  \item \textbf{Return every selected hit as a row.} A Boolean filter alone
    can keep or remove an input row. It cannot turn that row into two rows.
\end{itemize}

Figure~\ref{fig:scan-occurrences} uses abstract symbols to show how one input
row can produce two occurrences.

\begin{figure}[htbp]
\centering
\begin{tikzpicture}[font=\footnotesize]
  \node[anchor=west,font=\bfseries] at (0,3.15) {One stored row};
  \node[scan input,minimum width=1.35cm] at (.85,2.28) {id 27};
  \foreach \x/\s in {2.1/X,2.85/Y,3.6/X}
    \node[scan cell] at (\x,2.28) {\texttt{\s}};
  \foreach \x/\p in {2.1/0,2.85/1,3.6/2}
    \node[anchor=south,text=black!65] at (\x,2.6) {\p};
  \node[scan hit,minimum width=1.4cm] (query) at (2.85,.65)
    {query \texttt{X}};
  \draw[scan arrow,scanOrange] (query.north west) -- (2.1,1.99);
  \draw[scan arrow,scanOrange] (query.north east) -- (3.6,1.99);
  \node[anchor=west,font=\bfseries] at (7.0,3.15) {Two result rows};
  \node[scan output,minimum width=4.8cm] (r0) at (9.45,2.28)
    {id 27\qquad hit start 0};
  \node[scan output,minimum width=4.8cm] (r2) at (9.45,1.17)
    {id 27\qquad hit start 2};
  \draw[scan arrow] (4.15,2.28) -- (r0.west);
  \draw[scan arrow] (4.15,2.28) to[out=-35,in=180] (r2.west);
  \node[anchor=north,text=black!65] at (9.45,.60)
    {the identifier repeats, the hit changes};
\end{tikzpicture}
\caption{Searching \texttt{XYX} for \texttt{X} gives two occurrences.
The output must preserve both even though they come from the same stored row.}
\label{fig:scan-occurrences}
\end{figure}

For $N_{\mathrm{in}}$ input rows, if row $i$ produces $m_i$ hits, the number
of output rows is
\[
  N_{\mathrm{out}}=\sum_{i=1}^{N_{\mathrm{in}}}m_i,\qquad m_i\geq0.
\]
The repeated identifier is intentional. It does not mean that the table
contains duplicate records.

\subsection{Why not pass a query as text?}

An early interface idea was to pass a table name and an SQL string to a search
function. The following is a design sketch, not a BioSeqDB API:

\noindent\begin{minipage}{\linewidth}
\begin{lstlisting}[style=bioseq-search-sql,caption={An early interface sketch passes SQL as text.}]
SELECT * FROM search_sql(
  'targets',
  'SELECT id, seq FROM reads WHERE id > 20'
);
\end{lstlisting}
\end{minipage}

This hides choices that ordinary SQL makes explicit:

\begin{itemize}
  \item \textbf{Which output is the query sequence?} A function could demand
    a column called \texttt{query}, assume a fixed column position, or accept
    a separate column parameter. Renaming or reordering the inner output then
    becomes part of an extra API contract.
  \item \textbf{Where do surrounding conditions apply?} In a statement such
    as the sketch below, the outer planner cannot inspect a string as a
    relation. It cannot automatically insert the target identifier test into
    the search's input or choose a join order across the function boundary.
  \item \textbf{Who connects the two plans?} PostgreSQL can parse and optimize
    the inner SQL when the function submits it. The extension still has to
    interpret its output and arrange any rewrite that combines inner and
    outer conditions. Passing SQL text does not provide that integration.
\end{itemize}

\noindent\begin{minipage}{\linewidth}
\begin{lstlisting}[style=bioseq-search-sql,caption={An outer condition cannot automatically filter a hidden scan.}]
-- Hypothetical SQL-string interface
SELECT h.*
FROM search_sql('targets', 'SELECT id, seq FROM reads') AS h
WHERE h.target_id = 27;
-- The function must itself restrict the hidden target scan.
\end{lstlisting}
\end{minipage}

The function sees the target table through the string \texttt{'targets'}.
The outer condition cannot become a restriction on that hidden table scan
unless the interface explicitly passes it in or the extension adds a rewrite.
In the typed query above, PostgreSQL sees \texttt{t.id > 20} and the target
relation together.

\subsection{Why not save the last result in a function?}

Another idea is for a Boolean search function to save its result for later
score and CIGAR calls. This gives the
appearance of sharing work, but the SQL expression does not describe the
dependency. PostgreSQL does not guarantee left-to-right evaluation of Boolean
conditions. It may rearrange or skip parts of an expression.\footnote{PostgreSQL
17, \href{https://www.postgresql.org/docs/17/sql-expressions.html\#SYNTAX-EXPRESS-EVAL}
{Expression Evaluation Rules}.}

\noindent\begin{minipage}{\linewidth}
\begin{lstlisting}[style=bioseq-search-sql,caption={Hidden state does not establish an evaluation order.}]
-- Hypothetical functions with shared hidden state
SELECT id, last_score(seq)
FROM targets
WHERE save_search(seq, $1)
  AND last_score(seq) > 100;
\end{lstlisting}
\end{minipage}

The second condition cannot rely on the first having run. Declaring these
functions \texttt{VOLATILE} limits some planner transformations, but does not
establish a safe producer--consumer order between them. It also gives up
useful assumptions about repeatable calls.\footnote{PostgreSQL 17,
\href{https://www.postgresql.org/docs/17/xfunc-volatility.html}
{Function Volatility Categories}.}
Nested searches, a rescan for a new outer row, and early query termination
would all need explicit state management. A session variable or temporary
table merely moves that responsibility elsewhere. The Boolean call still
cannot emit multiple hit rows.

\subsection{A set-returning function is a valid alternative}

SQL already has a direct way to express one-to-many output. A set-returning
function in a \texttt{LATERAL} join could return a record for every
hit.\footcite{postgresql17_set_returning_functions} The \texttt{LATERAL} clause
lets that call refer to the current target row.\footnote{PostgreSQL 17,
\href{https://www.postgresql.org/docs/17/queries-table-expressions.html\#QUERIES-LATERAL}
{\texttt{LATERAL} Subqueries}.}
This alternative solves both the
row-count problem and the dependency between result fields:

\noindent\begin{minipage}{\linewidth}
\begin{lstlisting}[style=bioseq-search-sql,caption={A different API could return explicit hit records.}]
-- Alternative API sketch, not an existing BioSeqDB function
SELECT t.id, h.score, h.start, h.cigar
FROM targets AS t
CROSS JOIN LATERAL align_hits(t.seq, $1) AS h
WHERE t.id > 20;
\end{lstlisting}
\end{minipage}

The reason to develop custom scans is therefore an engineering choice. We
wanted the search to appear in ordinary \texttt{WHERE} and \texttt{JOIN ON}
clauses, while one executor node controls batches, hit fields, and sequence
lifetimes. A set-returning API would expose the hit record directly and would
need its own design for batching several input rows.
